\documentclass[12pt,a4paper]{article}
\usepackage[utf8]{inputenc}
\usepackage[margin=1in]{geometry}
\usepackage{amsmath, graphicx, booktabs, listings, hyperref}

\title{Sample Technical Report}
\author{Pavan A Gowda}
\date{\today}

\lstset{basicstyle=\ttfamily\small, frame=single, breaklines=true}

\begin{document}
\maketitle
\tableofcontents
\newpage

\section{Introduction}
This report demonstrates core \LaTeX{} features: sections, lists, math,
tables, code listings and cross-references (see Section~\ref{sec:math}).

\section{Mathematics}\label{sec:math}
Inline math: $a^2 + b^2 = c^2$. A numbered equation:
\begin{equation}
  \int_0^\infty e^{-x^2}\,dx = \frac{\sqrt{\pi}}{2}
  \label{eq:gauss}
\end{equation}
Equation~\eqref{eq:gauss} is the Gaussian integral.

\section{Lists and Tables}
\begin{itemize}
  \item Version control: Git
  \item Hosting: GitHub
\end{itemize}

\begin{table}[h]
  \centering
  \begin{tabular}{lcc}
    \toprule
    Tool & Type & Output \\
    \midrule
    LaTeX   & Typesetting & PDF  \\
    Doxygen & Code docs   & HTML \\
    \bottomrule
  \end{tabular}
  \caption{Tools covered in TW.}
  \label{tab:tools}
\end{table}
Table~\ref{tab:tools} summarizes the tools.

\section{Code}
\begin{lstlisting}[language=C]
int add(int a, int b) { return a + b; }
\end{lstlisting}

\section{Conclusion}
\LaTeX{} separates content from layout and produces consistent documents.

\end{document}
